\section{Fan holonomy and symplectic wall transvections}
\label{sec:holonomy}

The invariant-lattice description of \cref{thm:globalization-criterion} is
global.  This section turns it into an edge-local criterion and identifies the
full-rank wall maps as symplectic transvections.

\subsection{Rank-one relative transitions}

Let $R$ and $R'$ be adjacent maximal rigid charts.  Their presentation cones
share a codimension-one face.  Choose the primitive integral covector
\[
 \ell_{R\to R'}\in K^*
\]
which vanishes on the common wall and is positive on the new ray of $R'$.
The sign of $\ell$ is fixed by this convention, but the results below are
unchanged if both $v$ and $\ell$ are negated.

\begin{proposition}[Integral rank-one wall form]
\label{prop:rank-one-wall}
For adjacent charts $R,R'$, there is an integral column vector $u$ such that
\begin{equation}
 \mathsf L_{R'}-\mathsf L_R=u\ell^T.
 \label{eq:rank-one-difference}
\end{equation}
The relative transition
\begin{equation}
 T_{R\to R'}:=\mathsf L_R^{-1}\mathsf L_{R'}
 \label{eq:relative-transition}
\end{equation}
has the form
\begin{equation}
 \boxed{T_{R\to R'}=I+v\ell^T,\qquad v=\mathsf L_R^{-1}u\in K.}
 \label{eq:rank-one-transition}
\end{equation}
Moreover
\begin{equation}
 \ell^Tv=0,
 \qquad
 \det T_{R\to R'}=1.
 \label{eq:unipotent-wall}
\end{equation}
Thus $(T-I)^2=0$.
\end{proposition}

\begin{proof}
The two local maps $\mathsf L_R$ and $\mathsf L_{R'}$ agree on their common
wall, because a weight on that wall has the same rigid direct-sum
representative and the same cluster $g$-vector in both charts.  Hence
$\mathsf L_{R'}-\mathsf L_R$ vanishes on $\ker\ell$, so it has rank at most
one and factors as $u\ell^T$ over $\mathbb Q$.  Since the difference is
integral and $\ell$ is primitive, Bezout's identity for the entries of
$\ell$ implies $u\in\mathbb Z^n$.  Multiplying by the unimodular matrix
$\mathsf L_R^{-1}$ gives \cref{eq:rank-one-transition} with integral $v$.

Along one mutation, both $G_R$ and $\Delta_R$ are multiplied on the right by
unimodular mutation matrices of determinant $-1$.  Therefore
\[
 \det\mathsf L_R=\frac{\det G_R}{\det\Delta_R}
\]
is constant on the connected exchange graph.  Hence
$\det T_{R\to R'}=1$.  The matrix determinant lemma gives
\[
 \det(I+v\ell^T)=1+\ell^Tv,
\]
so $\ell^Tv=0$.  Finally
$(v\ell^T)^2=v(\ell^Tv)\ell^T=0$.
\end{proof}

\begin{remark}
The vector $v$ is the shear direction and $\ell$ is the wall normal.  The
condition $\ell^Tv=0$ says that the shear direction lies in the wall itself.
Zero-shear edges are precisely those for which $v=0$, equivalently
$\mathsf L_R=\mathsf L_{R'}$.
\end{remark}

\subsection{The symplectic criterion}

\begin{theorem}[Rank-one preservation criterion]
\label{thm:rank-one-symplectic}
Let $S\in\Skew_n(\mathbb Z)$ and let
\[
 T=I+v\ell^T\in\mathrm{GL}_n(\mathbb Z),
 \qquad \ell^Tv=0,
\]
with $\ell$ primitive and nonzero.  Then
\begin{equation}
 \boxed{
 T^TST=S
 \quad\Longleftrightarrow\quad
 Sv=c\ell\text{ for some }c\in\mathbb Z.}
 \label{eq:symplectic-transvection-criterion}
\end{equation}
If $S$ is nondegenerate and $v\ne0$, then $c\ne0$.
\end{theorem}

\begin{proof}
Since $S$ is skew-symmetric, $v^TSv=0$.  Expanding gives
\begin{align*}
 T^TST-S
 &=\ell v^TS+Sv\ell^T+\ell(v^TSv)\ell^T\\
 &=-\ell(Sv)^T+(Sv)\ell^T.
\end{align*}
This vanishes exactly when the two column vectors $Sv$ and $\ell$ are
proportional over $\mathbb Q$.  Because both are integral and $\ell$ is
primitive, the proportionality scalar is an integer.  If $S$ is
nondegenerate and $v\ne0$, then $Sv\ne0$, so the scalar is nonzero.
\end{proof}

\begin{definition}[Integral $S$-transvection]
A nonidentity matrix of the form
\[
 I+v\ell^T,
 \qquad \ell^Tv=0,
 \qquad Sv=c\ell,
\]
is called an \emph{integral $S$-transvection}.  When $S$ is nondegenerate,
this is a symplectic transvection for the integral symplectic lattice
$(K,S)$.
\end{definition}

\begin{corollary}[Edge-local globalization test]
\label{cor:edge-local-test}
A skew form $S$ is globalizable if and only if it is preserved by the relative
transition of every oriented edge in a connected spanning set of the
reachable exchange graph.  For a nonzero rank-one edge this is equivalent to
\[
 Sv_{R\to R'}\in\mathbb Z\ell_{R\to R'}.
\]
Thus $\Qglob$ is obtained by intersecting elementary wall conditions.
\end{corollary}

\begin{proof}
The relative edge transitions generate the transfer group $\Hol$.  Apply
\cref{thm:globalization-criterion,thm:rank-one-symplectic}.
\end{proof}

\subsection{Piecewise symplecticity in full rank}

\begin{theorem}[Symplectic presentation-to-cluster transport]
\label{thm:piecewise-symplectic}
Assume $B_0^T\Lambda_0=dI_n$.  Then:
\begin{enumerate}[label=(\roman*),leftmargin=2.4em]
\item every chamber transfer is integral symplectic,
      \begin{equation}
       \mathsf L_R^T\Lambda_0\mathsf L_R=\Lambda_0;
       \label{eq:piecewise-symplectic}
      \end{equation}
\item every nonzero adjacent relative transition is an integral
      $\Lambda_0$-transvection;
\item the fan map $\Gamma_T$ is piecewise integral symplectic, and its
      nonlinearity across a wall is completely determined by the pair
      $(v_{R\to R'},\ell_{R\to R'})$.
\end{enumerate}
\end{theorem}

\begin{proof}
Part (i) is \cref{cor:full-rank-globalizes} rewritten through
\cref{thm:globalization-criterion}.  Relative transitions of symplectic
matrices are symplectic.  Apply
\cref{prop:rank-one-wall,thm:rank-one-symplectic} to obtain (ii).  Part (iii)
is the chamberwise interpretation of (i) and (ii).
\end{proof}

\begin{remark}[Comparison with tropical duality]
Tropical duality relates the $G$- and $C$-matrices of one cluster pattern
\cite{NakanishiZelevinsky}.  The matrix $\mathsf L_R=G_R\Delta_R^{-1}$
compares two different but object-labelled patterns: the cluster $g$-chart and
the split presentation chart.  The symplectic theorem is therefore not a
restatement of $G_R^TC_R=I$; it uses the suspended companion seed and the
quantum compatibility form.
\end{remark}

\subsection{Holonomy, zero-shear components and invariant forms}

The edge description gives a useful geometric decomposition.  Remove every
nonzero-shear edge from the maximal-rigid exchange graph.  On each remaining
connected component, all matrices $\mathsf L_R$ are equal, so $\Gamma_T$ is
represented by one integral matrix.  We call these the \emph{linearity
components}.  Contracting each component yields a graph whose edges are
rank-one transvections.  The transfer group is generated by these contracted
edge maps.

\begin{proposition}[Invariant forms as holonomy fixed points]
\label{prop:holonomy-fixed}
Let $\mathcal G$ be the contracted graph of linearity components and choose a
spanning tree.  Then $\Qglob$ is the intersection of:
\begin{enumerate}[label=(\roman*),leftmargin=2.3em]
\item the fixed lattices of the transvections on the spanning tree; and
\item the fixed lattices of the cycle holonomies corresponding to the
      non-tree edges.
\end{enumerate}
In particular, if the transvections associated with a generating edge set
force a one-dimensional rational invariant space, then every globalizable
integral parameter is an integral multiple of one primitive skew form.
\end{proposition}

\begin{proof}
Choose a base component.  Products along the spanning tree recover the
transfer of every vertex from the base.  Each non-tree edge adds one generator
of the fundamental group and hence one cycle holonomy.  A form is fixed by all
vertex transfers exactly when it is fixed by this generating set.  The final
statement follows by intersecting the rational invariant line with the
integral skew lattice.
\end{proof}

\begin{algorithm}[Computing the globalization lattice]
\label{alg:globalization-lattice}
Given a finite paired presentation/cluster atlas:
\begin{enumerate}[label=\arabic*.,leftmargin=2.2em]
\item compute $\Delta_R,G_R$ and
      $\mathsf L_R=G_R\Delta_R^{-1}$ for every maximal rigid chart;
\item deduplicate equal transfer matrices and assemble the integer equations
      $\mathsf L_R^TS\mathsf L_R-S=0$ on the
      $\binom n2$ upper-triangular entries of $S$;
\item compute the rational nullspace and saturate it in
      $\Skew_n(\mathbb Z)$;
\item independently factor every adjacent transition as
      $I+v\ell^T$ and verify the equivalent wall equations
      $Sv\in\mathbb Z\ell$.
\end{enumerate}
The two computations provide independent certificates of the same lattice.
\end{algorithm}
